\chapter{Graphs, supports and physically admissible groups}\label{ch:graphs}
An edge is a possible physical transformation, not a promise that a helpful
transfer can always occur. A route, a pump and a warehouse doorway may each
impose limits independent of any potential. Begin with those limits, because
an exact value for an impossible action is only a counterfactual calculation.

\section{Three graphs with different jobs}
The physical graph lists allowed transformations. The potential-factor graph
lists which state coordinates a term in $V$ depends on. The action-interaction
graph records overlap or coupling between actions. In the diagonal Gaussian
model every potential factor has one coordinate, yet two physical actions
can interact because they both change the same source. A diagonal Hessian
does not make all edge actions independent.

For action $a$, its write support is the set of coordinates with nonzero
increment $\delta_a$. Its constraint support includes capacities, reservations
or aggregate budgets it uses. Two actions may have disjoint write supports
but share a machine capacity. Conversely, reading the same immutable scale
does not by itself create a shared physical constraint.

\section{Joint feasibility is not singleton feasibility}
Suppose Ridge contains five units. Two requests each withdraw three units
from Ridge to different destinations. Each is individually nonnegative at
its endpoint, but their group withdraws six and is invalid. Testing children
one at a time against the same five units does not validate their sum.

With $q_a\geq0$, a conservative frozen-stock rule might require
\begin{equation}
 \sum_{a:\,\mathrm{src}(a)=i}q_a\leq z_i.
\end{equation}
This is an example of a declared rule, not a newly adopted programme policy.
A plant that permits concurrent incoming flow to fund outgoing flow requires
a different timing and feasibility contract. Endpoint nonnegativity alone
does not prove that an unmodelled intermediate mechanism can perform it.

Destination capacity is another independent constraint. If a tank can hold
ten and begins at eight, two incoming deliveries of two each cannot both
fit merely because their separate endpoints do. Physical capacity is not
the Gaussian reference. A reference may be six even when the tank can hold
ten; exceeding the reference can be feasible but increase potential.

\section{The straight segment and its limited guarantee}
The simplex is convex. If $z$ and $z+\delta_G$ lie in it, then
\begin{equation}
 z+s\delta_G=(1-s)z+s(z+\delta_G),\qquad 0\leq s\leq1,
\end{equation}
also lies in it. This proves nonnegative conserved stock along the aggregate
straight interpolation. It does not prove child-specific rate limits,
mechanical synchronization, queue clearance or causal access to incoming
stock. Those are additional constraints on a physical realization.

This distinction matters when a path integral is later described as physical.
The same segment can be a valid mathematical interpolation even when the
plant has not declared constant-rate joint execution. The integral identity
survives; the physical-path interpretation does not follow automatically.

\section{Grouping without losing atomic actions}
A joint group retains child identities and increments. It is not an excuse
to conceal a fuel-to-electricity-to-pump-to-water chain inside one undifferentiated
operation. Distinct carriers and intermediate transformations must remain
traceable when they belong to the represented scope.

The analytical bridge \cite{bridge} groups coupled actions under a common
accounting boundary and distinguishes physical grouping from receipt batching.
Batching independent records can reduce administrative overhead, but it does
not create a new physical synergy. Conversely, simultaneous actions sharing
a source need joint evaluation even if their paperwork is submitted separately.

\section*{Exercises with reasoning}
Take a source with six units and two proposed quantities $(4,4)$. A common
cap of six with proportional scaling would yield $(3,3)$. The number three
comes from that resolver, not from the value equation. If the resolver were
priority-based, a different accepted vector could be physically admissible.
Never price $(4,4)$ and execute $(3,3)$ under the old value record.

Now let two accepted actions use disjoint nodes and disjoint constraints.
With a separable potential, their finite values add because the affected
terms are disjoint. If a new pair factor couples one node from each action,
the independence proof no longer applies even though the physical edges
remain disjoint. The correct affected support follows the potential factors.

\section*{Chapter recap}
The graph says what can connect, support says what changes, constraints say
what is permitted, and the potential says how the represented change is
valued. None can silently perform the job of another.
